\section*{Chapter \ref{ch:m3:crude-oil} --- Crude Oil}

\begin{solution}{exo:m3:crude-oil:1}
$141.5/0.87 - 131.5 = 31.1$° API: heavier (denser) than a 39.6° crude.
\end{solution}

\begin{solution}{exo:m3:crude-oil:2}
The 25th of November 2026 is a business day, so trading ends three business days
before it: Friday 20 November 2026.
\end{solution}

\begin{solution}{exo:m3:crude-oil:3}
A single field's output falls over time and a small deliverable supply can be
squeezed: whoever controls most cargoes of a month can move the price. Several grades,
with quality adjustments and the most-competitive-grade rule, make the deliverable
supply larger and the benchmark harder to manipulate.
\end{solution}

\begin{solution}{exo:m3:crude-oil:4}
$80 + 0.04 \times (36 - 38) - 0.12 \times (0.60 - 0.40)/0.1 = 80 - 0.08 - 0.24 =
\qty{79.68}{\$/\barrel}$.
\end{solution}

\begin{solution}{exo:m3:crude-oil:5}
$25 - 0.40 = \qty{24.60}{\$/\barrel}$. A front price of \$18 implies a storage cost of
\$7 per barrel-month: storage at \$0.40 is not available to those who could arbitrage,
or the front contract is being liquidated faster than storage can be found.
\end{solution}

\begin{solution}{exo:m3:crude-oil:6}
$(71.35 + 2.20) \times 2\,000\,000 = $ USD~147.1 million.
\end{solution}

\begin{solution}{exo:m3:crude-oil:7}
25 days above \$5; the widest spread is \$58.06 (20 April 2020), the second-widest
\$8.49.
\end{solution}

\begin{solution}{exo:m3:crude-oil:8}
A buyer of a physically delivered contract takes oil at a specific place and date.
If there is nowhere to put it and moving it costs money, the oil at that place and
date has a negative value to anyone without storage: it is a liability to be
disposed of. ``Free oil'' is worth its value elsewhere, later, minus the cost of
storing and moving it, which can exceed that value.
\end{solution}

\begin{solution}{pb:m3:crude-oil:1}
\textbf{1.} 1\,000\,000 barrels. \textbf{2.} 21 April 2020: 25 April was a
Saturday, the last business day before it Friday 24 April, and three business days
before that is Tuesday 21 April; 20 April is the day before. \textbf{3.} USD~18.27
million. \textbf{4.} \$6.76, a storage cost of \$6.76 per barrel-month. \textbf{5.}
No: about 17 times the normal lease; storage was already scarce.
\textbf{6.} $(-37.63 - 18.27) \times 1\,000\,000 = -$USD~55.90 million.
\textbf{7.} \$58.06 per barrel-month. \textbf{8.} About 145 times. \textbf{9.}
Nobody who could still take delivery and store had space left, so the arbitrage of
\cref{prop:m3:crude-oil:floor} could not be done at scale. \textbf{10.}
$(58.06 - 0.40) \times 1\,000\,000 = $ USD~57.66 million.
\textbf{11.} \$6.76 a barrel, USD~6.76 million. \textbf{12.} \$6.76 a barrel, 27\%
of the June price, in a month. \textbf{13.} Later months were further from delivery
and priced with the expected recovery; a fund in them was not forced to sell into the
settlement window. \textbf{14.} The price was set by the average of spread trades
between 14:28 and 14:30, when holders without storage were forced to sell: a small
window at the worst moment. \textbf{15.} The Bachelier model, which allows a negative
underlying; lognormal models cannot price options on a negative future.
\textbf{16.} Holders of storage or pipeline space, refiners, and traders short the
May contract with delivery capacity: for them oil at $-\$37.63$ was paid storage.
\textbf{17.} A physically settled benchmark prices the commodity at one place and
date, including the cost of taking it; a cash-settled one prices an index and cannot
be forced into delivery. \textbf{18.} No: ICE Brent is cash-settled on an index of a
seaborne market with ships as storage; the EIA's Brent spot price fell no lower than
\$9.12 (21 April 2020). \textbf{19.} \emph{Named result:} \$58.06 per barrel for one
month, about 145 times the illustrative lease. \textbf{20.} Of the commodity at the
delivery place and date, net of what it costs to take and keep it there.
\end{solution}

\begin{solution}{iq:m3:crude-oil:1}
\omterm{def:m3:crude-oil:brent}{Dated Brent} is an agency's assessment of physical North Sea cargoes loading in the
coming weeks; ICE Brent futures are exchange-traded contracts for a month months
ahead, cash-settled against an index of the forward market. The CFDs and the forward
market link them.
\iqlookfor{physical versus paper, the time dimension, the linking instruments.}
\end{solution}

\begin{solution}{iq:m3:crude-oil:2}
When the price is for delivery at a place and time where taking the commodity costs
more than it is worth: full storage, no transport, disposal costs (power in a
congested hour, gas in a full network). Paper follows the physical at delivery.
\iqlookfor{the price is of a delivered good; storage and logistics set a floor only if available.}
\end{solution}

\begin{solution}{iq:m3:crude-oil:3}
Whether the continuous series rolled before the penultimate day as any real fund
would, whether the strategy could have held the May contract at all (delivery,
accountability levels), whether the fill at the settlement was achievable, and
whether the P\&L comes from one day; a log-return series breaks on a negative price.
\iqlookfor{roll conventions, feasibility, the negative price as a data bug trap.}
\end{solution}

\begin{solution}{iq:m3:crude-oil:4}
A spread of \$1.50 against \$0.40 of storage pays \$1.10 for cash-and-carry: buy
the front, take delivery, store, sell the second month. Risks: storage and pipeline
scheduling, financing, quality and losses, and the possibility that the spread widens
further before delivery, which costs margin.
\iqlookfor{the carry trade, and its operational and liquidity risks.}
\end{solution}

\begin{solution}{iq:m3:crude-oil:5}
$\ln F$ is undefined for $F \le 0$: prices, implied volatilities and Greeks fail.
Switch to a normal (Bachelier) model or a shifted lognormal, re-mark the surface in
that model, and make the code refuse a log of a non-positive price rather than return
NaN.
\iqlookfor{Bachelier, shifted models, defensive code.}
\end{solution}

\begin{solution}{iq:m3:crude-oil:6}
Several grades mean more deliverable supply and more sellers, so cornering the
cheapest grade is expensive; the price is the most competitive grade's. The cost is
complexity: quality premiums must be set and revised, and the benchmark's grade can
change from day to day, which adds basis for hedgers.
\iqlookfor{deliverable supply against squeezes; quality-adjustment and basis costs.}
\end{solution}
